\documentclass[11pt,a4paper,oneside]{book}

% --- Paquetes Básicos ---
\usepackage[utf8]{inputenc}
\usepackage[spanish,es-tabla]{babel}
\usepackage[T1]{fontenc}
\usepackage{graphicx}
\usepackage{geometry}
\usepackage{hyperref}
\usepackage{xurl}
\usepackage{titlesec}
\usepackage{booktabs}
\usepackage{textcomp}
\usepackage{eurosym}
\usepackage{array}
\usepackage{caption}
\usepackage{float}
\usepackage{microtype}
\usepackage{amsmath}
\usepackage{amssymb}
\usepackage{tikz}
\usetikzlibrary{babel, positioning, shapes.geometric}

% --- Configuración de Página ---
\geometry{
    a4paper,
    top=3cm,
    bottom=3cm,
    left=3cm,
    right=3cm
}

% --- Datos del TFG ---
\title{Sistema de Predicción de Demanda para Retail con Análisis de Incertidumbre\\[0.5em]
\large Borrador para revisión --- Capítulos 1--3}
\author{Artem Mindlin}
\date{\today}

\begin{document}

\maketitle

\frontmatter
% !TEX root = ../main.tex
\chapter*{Resumen}
\addcontentsline{toc}{chapter}{Resumen}

El sector retail se enfrenta al reto constante de equilibrar la disponibilidad de producto con la minimización de mermas. Este Trabajo de Fin de Grado presenta el desarrollo de un Sistema de Apoyo a la Decisión (DSS) End-to-End cuyo objetivo operativo es estimar y automatizar cuántas unidades debe pedir diariamente cada SKU bajo restricciones logísticas reales.

La metodología propuesta evoluciona el \textit{forecasting} clásico integrando modelos de Gradient Boosting con \textit{Mondrian Conformal Prediction} para garantizar la equidad estadística de la incertidumbre por categoría de producto. El sistema no se limita a predecir, sino que transforma estos pronósticos probabilísticos en decisiones de inventario mediante la formulación del \textit{Newsvendor}, que traduce la distribución de demanda estimada en una cantidad de pedido según la asimetría entre el coste de rotura y el de sobrestock. Finalmente, el modelo estadístico se ha operativizado bajo estándares MLOps, incluyendo explicabilidad automática (SHAP), trazabilidad de experimentos por artefactos versionados y un cuadro de mandos interactivo renderizado en servidor.

Los resultados empíricos demuestran que la corrección de la demanda censurada, bajo una evaluación justa que enfrenta a todas las estrategias contra una misma demanda real de referencia, reduce el coste logístico de inventario en un \pendiente{92{,}4}\% (con un incremento en el fill rate del \pendiente{60{,}2}\% al \pendiente{97{,}6}\%), y que, a coste de inventario prácticamente idéntico, el modelo probabilístico CatBoost eleva el nivel de servicio del \textbf{79,3\%} al \textbf{94,4\%} frente al modelo estacional ingenuo, pese a que este último obtiene el mejor MAE. Este enfoque integral contribuye directamente al \textbf{ODS 12 (Producción y Consumo Responsables)} al mitigar el desperdicio de productos perecederos mediante una cadena de suministro altamente eficiente, matemática e inteligible.

\textbf{Palabras clave:} Retail, Machine Learning, Investigación de Operaciones, Mondrian Conformal Prediction, Newsvendor, MLOps, ODS 12.

\newpage

\begin{otherlanguage}{catalan}
\chapter*{Resum}
\addcontentsline{toc}{chapter}{Resum}

El sector del retail s'enfronta al repte constant d'equilibrar la disponibilitat de producte amb la minimització de minves. Aquest Treball de Fi de Grau presenta el desenvolupament d'un Sistema de Suport a la Decisió (DSS) End-to-End amb l'objectiu operatiu d'estimar i automatitzar quantes unitats s'han de demanar diàriament per a cada SKU sota restriccions logístiques reals.

La metodologia proposada evoluciona el \textit{forecasting} clàssic integrant models de Gradient Boosting amb \textit{Mondrian Conformal Prediction} per a garantir l'equitat estadística de la incertesa per categoria de producte. El sistema no es limita a predir, sinó que transforma aquests pronòstics probabilístics en decisions d'inventari mitjançant la formulació del \textit{Newsvendor}, que tradueix la distribució de demanda estimada en una quantitat de comanda segons l'asimetria entre el cost de ruptura i el de sobreestoc. Finalment, el model estadístic s'ha operativitzat sota estàndards MLOps, incloent-hi explicabilitat automàtica (SHAP), traçabilitat d'experiments per artefactes versionats i un quadre de comandament interactiu renderitzat al servidor.

Els resultats empírics demostren que la correcció de la demanda censurada, sota una avaluació justa que enfronta totes les estratègies contra una mateixa demanda real de referència, redueix el cost logístic d'inventari en un \pendiente{92{,}4}\% (amb un increment del fill rate del \pendiente{60{,}2}\% al \pendiente{97{,}6}\%), i que, a un cost d'inventari pràcticament idèntic, el model probabilístic CatBoost eleva el nivell de servei del \textbf{79,3\%} al \textbf{94,4\%} respecte del model estacional ingenu, tot i que aquest darrer obté el millor MAE. Aquest enfocament integral contribueix directament a l'\textbf{ODS 12 (Producció i Consum Responsables)} en mitigar el malbaratament de productes peribles mitjançant una cadena de subministrament altament eficient, matemàtica i intel·ligible.

\textbf{Paraules clau:} Retail, Machine Learning, Investigació d'Operacions, Mondrian Conformal Prediction, Newsvendor, MLOps, ODS 12.
\end{otherlanguage}

\newpage

\begin{otherlanguage}{english}
\chapter*{Abstract}
\addcontentsline{toc}{chapter}{Abstract}

The retail sector faces the constant challenge of balancing product availability with waste minimization. This Bachelor's Thesis presents the development of an End-to-End Decision Support System (DSS) whose operational goal is to estimate and automate how many units of each SKU should be ordered daily under real logistical constraints.

The proposed methodology evolves classical forecasting by integrating Gradient Boosting models with Mondrian Conformal Prediction to ensure statistical fairness of uncertainty across product categories. The system goes beyond prediction by transforming probabilistic forecasts into inventory decisions via the \textit{Newsvendor} formulation, which turns the estimated demand distribution into an order quantity according to the asymmetry between stockout and overstock costs. Finally, the statistical model has been operationalized under MLOps standards, including automated explainability (SHAP), experiment traceability through versioned run artifacts, and an interactive server-rendered dashboard.

Empirical results demonstrate that correcting censored demand, under a fair evaluation that scores every strategy against one and the same real reference demand, reduces inventory logistics costs by \pendiente{92.4}\% (improving the fill rate from \pendiente{60.2}\% to \pendiente{97.6}\%), and that, at practically identical inventory cost, the probabilistic CatBoost model raises the service level from \textbf{79.3\%} to \textbf{94.4\%} against the seasonal naïve baseline, even though the latter attains the best MAE. This comprehensive approach directly contributes to \textbf{SDG 12 (Responsible Consumption and Production)} by mitigating perishable product waste through a highly efficient, mathematical, and intelligible supply chain.

\textbf{Keywords:} Retail, Machine Learning, Operations Research, Mondrian Conformal Prediction, Newsvendor, MLOps, SDG 12.
\end{otherlanguage}

\tableofcontents
\listoffigures
\listoftables

\mainmatter

% !TEX root = ../main.tex
\chapter{Introducción}
\label{chap:introduccion}

Cada mañana, en cada tienda de una cadena de distribución, alguien decide cuántas unidades de cada
producto fresco va a pedir. Lo hace sin conocer la demanda que tendrá que atender, y equivocarse
cuesta dinero en las dos direcciones: si pide de menos pierde la venta y el cliente se encuentra el
hueco en el lineal; si pide de más, el producto caduca y termina destruido. Esa decisión, repetida
miles de veces al día sobre artículos de vida útil corta y demanda volátil, es el problema que
aborda este trabajo. La gestión de la cadena de suministro en el sector \textit{retail} ha
evolucionado de un enfoque reactivo a uno proactivo apoyado en la analítica avanzada, pero la
previsión de la demanda de frescos sigue siendo uno de sus desafíos abiertos.

\section{Motivación}

\subsection{Motivación profesional}
La relevancia del problema no es solo operativa, sino también económica y ética. A nivel europeo, el
desperdicio alimentario asciende a unos 59 millones de toneladas anuales, con un coste estimado de
132.000 millones de euros para el conjunto de la cadena alimentaria
\cite{eucommission2023foodwaste}. En España, estas cifras han servido de base para una
transformación del marco legislativo con la Ley 1/2025 de Prevención de las Pérdidas y el
Desperdicio Alimentario \cite{leydesperdicio2025}, cuyas obligaciones para los agentes de la cadena
alimentaria rigen desde el 2 de abril de 2026. Entre ellas figura disponer de un plan de prevención
que concrete cómo se aplicará la jerarquía de usos de los excedentes, y su incumplimiento se
sanciona con multas de hasta 500.000 euros en las infracciones muy graves. Quedan excluidas las
microempresas y los establecimientos de menos de 1.300 metros cuadrados, de modo que la obligación
recae precisamente sobre las grandes cadenas de distribución.

En este contexto, si bien la legislación regula qué hacer con el producto sobrante, priorizando la
donación para consumo humano antes que la alimentación animal, el aprovechamiento industrial o el
compostaje, la gestión logística de esos excedentes sigue suponiendo un coste para la empresa. Por
tanto, la herramienta tecnológica más eficaz para cumplir con la ley y proteger al mismo tiempo el
margen de beneficio es la prevención en origen, que es el propósito exacto del Sistema de Apoyo a la
Decisión desarrollado en este proyecto: evitar el exceso de inventario antes de que llegue a
producirse.

Este Trabajo de Fin de Grado nace de la necesidad de alinear los modelos matemáticos de previsión
con la realidad de la toma de decisiones en la reposición. El caso de uso central no es una
predicción abstracta de demanda, sino la construcción de un \textbf{Sistema de Apoyo a la Decisión
(DSS)} que responda a una pregunta operativa concreta del responsable de aprovisionamiento:
\emph{\textquestiondown cuántas unidades debe pedir hoy cada SKU para cubrir la demanda de los
próximos días?} Para contestarla, el sistema no solo estima la demanda esperada y cuantifica
matemáticamente su incertidumbre, sino que la traduce en una cantidad de pedido que pondera de forma
explícita el coste por rotura de stock (\textit{stockout}) frente al coste por sobrestock y
desperdicio, con un umbral de decisión (el fractil crítico) común al catálogo y una
curva de cuantiles propia de cada producto, contribuyendo así al Objetivo de Desarrollo Sostenible (ODS) 12: Producción y Consumo
Responsables \cite{ods12un}.

\subsection{Motivación personal}
El origen de este trabajo es anterior al propio Trabajo de Fin de Grado. Surgió mientras preparaba mi
candidatura a unas prácticas en el área de tecnología de una cadena de distribución alimentaria, para
la que quise construir por
mi cuenta un prototipo de previsión de demanda que sirviese como muestra de trabajo. Aquel prototipo
nunca llegó a terminarse. Resolvía la parte cómoda del problema, entrenar un modelo y mirar su
error, pero se detenía justo antes de las preguntas incómodas, que son las que ocupan esta memoria.
El TFG ha sido la ocasión de retomarlo, completarlo y someterlo al rigor experimental del que
entonces carecía.

La orientación del trabajo responde a una convicción que se ha ido formando durante su desarrollo: un
modelo solo genera valor cuando alguien puede usarlo para decidir. No se trata de restar importancia
al fundamento teórico, que ocupa el núcleo metodológico de esta memoria, sino de exigirle que llegue
hasta una decisión ejecutable en lugar de agotarse en la evaluación. Un cuaderno de análisis que
produce una cifra correcta no es todavía una herramienta, y mi impresión es que esa distancia
entre el resultado experimental y el sistema en funcionamiento es precisamente la que más
interesa resolver en el entorno empresarial. Por eso el proyecto no se detiene en la tabla de
métricas y llega hasta la API, el cuadro de mandos y el despliegue, con la trazabilidad necesaria
para que cualquier número publicado pueda rehacerse.

Finalmente no me incorporé al sector de la distribución alimentaria, sino a una empresa de
automoción, donde he encontrado un problema de la misma familia planteado sobre otro catálogo. Allí
los vehículos se piden al fabricante con mucha antelación respecto a la venta y, cuando no encuentran
comprador, se traspasan al canal de ocasión como kilómetro cero. La estructura de la decisión es la
misma que se estudia en esta memoria: cuánto pedir hoy frente a una demanda que no se conocerá
hasta después, con un coste del exceso distinto del coste del defecto. Lo que cambia es la
naturaleza de esa asimetría. El excedente no se destruye, se degrada de canal, de modo que el coste
del sobrestock es una pérdida de margen y no una merma total, sobre un catálogo de pocas unidades de valor alto en
lugar de muchas de valor bajo. Esa coincidencia estructural ha reforzado el enfoque adoptado aquí: lo
que se transfiere entre sectores no es el modelo concreto ni la parametrización de sus costes, sino
la forma de encadenar la estimación de la demanda, la cuantificación de su incertidumbre y la
decisión de aprovisionamiento que de ambas se deriva.

\section{Definición del Problema}
El enfoque tradicional del \textit{demand forecasting} se ha centrado históricamente en la obtención de predicciones puntuales (\textit{point forecasts}) mediante la minimización de métricas de error promedio como el MAE (\textit{Mean Absolute Error}) o el RMSE (\textit{Root Mean Squared Error}). Aunque útiles, estos enfoques presentan severas limitaciones cuando intentan utilizarse aisladamente en entornos de \textit{retail} de producción:

\begin{enumerate}
    \item \textbf{Censura de datos por stockout:} Cuando un producto se agota en el estante, la demanda observada se detiene, pero la demanda latente (lo que los clientes habrían comprado) sigue existiendo. Musalem et al.~\cite{musalem2010structural} demuestran empíricamente que ignorar este fenómeno sesga a la baja la demanda estimada de los artículos afectados. El sesgo, además, es doble: como establecen Anupindi et al.~\cite{anupindi1998stockout}, las ventas de los productos en rotura quedan censuradas por la derecha, mientras que las de los productos disponibles quedan infladas por las sustituciones que absorben esa demanda desviada, de modo que ninguna de las dos series mide ya la demanda real. Si este efecto no se corrige, el sistema de reposición perpetuará un ciclo de ineficiencia.
    \item \textbf{Desconexión entre error estadístico y coste logístico:} En productos frescos, el coste de equivocarse ``por exceso'' (desperdicio) es estructuralmente distinto al coste de equivocarse ``por defecto'' (pérdida de venta). Ban y Rudin~\cite{ban2019big} demuestran en el contexto del problema del vendedor de periódicos (\textit{Newsvendor}) que la precisión predictiva y el coste de inventario pueden divergir: un modelo puede mejorar su error un 20\% sin mejorar el coste, o incluso empeorándolo, cuando la estructura de costes es marcadamente asimétrica. Su propuesta es resolver el problema en un solo paso, minimizando empíricamente el coste del \textit{Newsvendor}, en lugar de estimar primero la demanda y optimizar después. Este trabajo recorre deliberadamente la vía en dos etapas, y la pieza que la sostiene es la función de pérdida asimétrica o \textit{pinball loss}. Un modelo entrenado con el error cuadrático aprende a predecir la demanda \emph{media}; entrenado con la \textit{pinball loss} aprende a predecir un \emph{cuantil} concreto de la distribución de demanda, que es precisamente lo que la decisión de pedido necesita. Y existe una equivalencia analítica exacta entre el fractil crítico del \textit{Newsvendor} y el cuantil que esa pérdida minimiza, resultado que se remonta a la regresión cuantílica introducida por Koenker y Bassett~\cite{koenkerbassett1978} y sistematizada posteriormente por Koenker~\cite{koenker2005quantile}. Sobre la rejilla de cuantiles estimada, la capa de decisión localiza después el fractil que corresponde a cada producto. Esta separación tiene una ventaja operativa que el enfoque en un solo paso no ofrece: la estructura de costes puede cambiar, o ser distinta para cada producto, sin necesidad de reentrenar el modelo.
\end{enumerate}

Este trabajo aborda ambas limitaciones mediante un enfoque unificado: desde un \textbf{forecasting probabilístico} avanzado con garantías matemáticas (\textit{Mondrian Conformal Prediction}) hasta la traducción de esa distribución en una decisión de pedido mediante Investigación de Operaciones (modelo del \textit{Newsvendor}), todo ello empaquetado en una arquitectura MLOps lista para despliegue.

\section{Objetivos}
\label{sec:objetivos}
El objetivo general de este proyecto es diseñar, evaluar y desplegar un sistema completo de apoyo a la reposición para \textit{retail} fresco, elevando la solución desde un experimento estadístico hasta un sistema desplegable que cierra el bucle entre predicción y decisión de reposición, con el fin último de minimizar el desperdicio alimentario y maximizar el margen de beneficio operativo.

Para alcanzar este objetivo general, se plantean los siguientes objetivos específicos:
\begin{enumerate}[label=\textbf{O\arabic*.}, ref=O\arabic*]
    \item \label{obj:pipeline} Desarrollar un \textit{pipeline} de datos reproducible que incluya la imputación de demanda censurada a partir del dataset \textit{Fresh RetailNet}, verificando que ninguna observación censurada por stockout introduce fuga de información hacia la etiqueta objetivo.
    \item \label{obj:modelos} Implementar modelos globales de \textit{Machine Learning} probabilístico (CatBoost, LightGBM) capaces de compartir información entre productos, obteniendo un coste de inventario simulado en el conjunto de test inferior al de un \textit{baseline} estacional ingenuo (\textit{seasonal naïve}).
    \item \label{obj:conformal} Garantizar la calibración estadística de las predicciones mediante \textit{Mondrian Conformal Prediction} con optimización multi-objetivo (\textit{Pareto Frontier}), alcanzando una cobertura empírica $\geq 80\%$ para el intervalo $[q_{0{,}1}, q_{0{,}9}]$ (nivel $\alpha = 0{,}20$) sobre el conjunto de test.
    \item \label{obj:mlops} Operativizar el sistema mediante estándares \textit{MLOps} e Ingeniería de Software,
    incluyendo trazabilidad de experimentos en un almacén de corridas, explicabilidad (SHAP),
    una API operacional y un cuadro de mandos interactivo \textit{What-If} renderizado en servidor
    (Django), soportado por \textit{Docker} e Integración Continua (CI/CD).
\end{enumerate}

\section{Alcance}
El presente trabajo se enmarca en el contexto de la previsión de demanda a corto plazo para productos frescos en el sector \textit{retail}. A continuación se delimitan explícitamente los límites del sistema para evitar interpretaciones erróneas sobre su cobertura.

\textbf{Dentro del alcance:}
\begin{itemize}
    \item Previsión de la demanda acumulada a nivel SKU-tienda para un horizonte de planificación de corto plazo (días), alineado con el ciclo de reposición habitual en distribución de frescos.
    \item Datos de granularidad diaria provenientes del benchmark \textit{Fresh RetailNet}, que cubre 50.000 combinaciones tienda-producto a lo largo de 90 días de observación.
    \item Corrección del sesgo por demanda censurada (stockout) mediante técnicas de imputación supervisada.
    \item Cuantificación de la incertidumbre con garantías estadísticas formales (Mondrian Conformal Prediction) y optimización multi-objetivo de hiperparámetros.
    \item Traducción de la distribución de demanda predicha en una cantidad de pedido por producto y día, mediante el fractil crítico del \textit{Newsvendor}, y evaluación del coste logístico y del nivel de servicio que esa política induce.
    \item Operativización del sistema mediante estándares MLOps (trazabilidad de ejecuciones, API operacional, CI/CD) y visualización interactiva.
\end{itemize}

\textbf{Fuera del alcance:}
\begin{itemize}
    \item \textbf{Optimización de precios}: los precios y descuentos se tratan como variables exógenas observables, no como variables de decisión del sistema.
    \item \textbf{Cadena de suministro multi-escalón}: el sistema opera a nivel de tienda individual; la gestión de almacenes centrales, proveedores o logística de transporte queda excluida.
    \item \textbf{Predicción a medio y largo plazo}: no se aborda la planificación estacional ni la previsión más allá del horizonte de reposición inmediato.
    \item \textbf{Simulación de inventario con estado}: la evaluación económica resuelve una
    decisión independiente por origen de predicción. No se modelan el arrastre de existencias entre
    periodos, los pedidos en tránsito ni el \textit{lead time}, por lo que el sistema ordena
    políticas pero no reproduce un libro de reposición. Tampoco se aborda el reparto de una
    capacidad de almacén limitada entre productos cuando la suma de las cantidades individualmente
    óptimas la supera. Ambas extensiones se detallan en la Sección~\ref{sec:trabajo_futuro}.
    \item \textbf{Ejecución automática de pedidos}: el sistema genera recomendaciones de pedido cuantificadas, pero la decisión final y la emisión del pedido corresponden al responsable de aprovisionamiento.
    \item \textbf{Productos no perecederos}: la metodología está orientada a frescos con alta rotación y corta vida útil; su extrapolación a otras categorías requeriría validación adicional.
\end{itemize}

\section{Impacto Esperado}
El sistema no se dirige a un único perfil de usuario. Las decisiones de reposición implican a
personas con responsabilidades distintas, y el beneficio que cada una obtiene es también distinto:

\begin{itemize}
    \item \textbf{Responsable de aprovisionamiento} (usuario operativo principal). Recibe, para cada
    producto y cada día, una cantidad de pedido recomendada acompañada del intervalo de demanda que
    la justifica, en lugar de una cifra aislada sin indicación de su fiabilidad. Dispone además de un
    simulador que permite contrastar el efecto de modificar los supuestos de coste antes de
    comprometer el pedido, de modo que la herramienta apoya la decisión sin sustituirla.
    \item \textbf{Responsable de tienda u operaciones} (lectura agregada). Dispone del nivel de
    servicio y del coste agregados de la tienda, y de un listado de excepciones que señala los
    productos en riesgo, de modo que puede localizar dónde se concentran las roturas y las mermas
    sin revisar el catálogo entero. La agregación por categoría de producto queda como extensión
    natural: la categoría ya estructura la calibración conformal, pero el sistema no publica hoy
    coste ni nivel de servicio agrupados por ella.
    \item \textbf{Equipo técnico responsable del sistema}. Cada ejecución queda registrada con su
    configuración, sus métricas y la versión exacta del código que la produjo, y el sistema
    contempla el reentrenamiento periódico y la vigilancia de la degradación del modelo. El efecto
    esperado es reducir el coste de mantener el sistema en el tiempo, que es donde suelen fracasar
    los proyectos de analítica que sí funcionaron en su evaluación inicial.
    \item \textbf{Comunidad académica y docente}. La implementación completa, sobre un conjunto de
    datos público y con los experimentos reproducibles, sirve como referencia para un problema que
    la literatura suele abordar por partes separadas. El detalle de este legado se recoge en la
    Sección~\ref{sec:legado}.
\end{itemize}

Para la organización en su conjunto, el impacto se mide en las dos magnitudes que la operación
percibe: el coste logístico total y el nivel de servicio. Los resultados del
Capítulo~\ref{chap:resultados} cuantifican ambas mejoras y permiten estimar el efecto de trasladar
la metodología a un catálogo real. Ese mismo efecto tiene una lectura ambiental y regulatoria
directa, ya que el exceso de inventario en productos perecederos termina en merma: reducirlo en
origen es la vía más eficaz de contribuir al ODS~12 y de sustentar el Plan de Prevención que exige
la normativa vigente. El análisis detallado de la vinculación con los Objetivos de Desarrollo
Sostenible se recoge en el Anexo~\ref{anx:ods}.

\section{Aportaciones del Trabajo}
\label{sec:aportaciones}
Este trabajo no propone una técnica nueva, sino que establece cómo debe evaluarse la
recuperación de demanda censurada, y lo hace sobre un \textit{benchmark} publicado en mayo de
2025~\cite{freshretailnet2025}. Sus aportaciones son las siguientes.

\begin{enumerate}
    \item \textbf{El protocolo estándar para medir la reconstrucción de demanda censurada es
    circular, y ninguna métrica calculada sobre él ordena reglas de reconciliación.} El protocolo
    fabrica su verdad-terreno reduciendo la venta de días limpios en proporción a una rotura
    inventada, de modo que una regla que se limite a deshacer esa proporción recupera la
    respuesta exacta sin modelo alguno y obtiene el menor error de reconstrucción de todas las
    estrategias evaluadas, siendo la que menos garantía ofrece sobre una rotura real. El
    favoritismo alcanza igualmente a la evaluación por coste, que se puntúa sobre los mismos días.
    De ahí que la elección entre reglas se argumente aquí sobre bases de modelado, que la
    comparación de hiperparámetros dentro de una misma regla siga siendo válida, y que la ventaja
    medida para la corrección de censura se declare como cota superior
    (Sección~\ref{sec:imputacion_censura_sintetica}).

    \item \textbf{Una partición de validación puede decidir la respuesta de una búsqueda de
    hiperparámetros, y hacerlo en silencio.} El imputador reajusta sobre el panel que recibe, de
    modo que partir ese panel para obtener un conjunto de validación disjunto encoge también el
    conjunto de entrenamiento del maestro. Medida sobre un maestro así empobrecido, la
    sintonización parecía mejorar un 5{,}33\%; los mismos hiperparámetros resultaban un 12{,}44\%
    \emph{peores} que no sintonizar a la escala real de operación, de modo que el fichero
    persistido era activamente perjudicial allí donde se consumía. Separar por tiempo y aplicar la
    partición como una máscara sobre el panel completo, en lugar de como un recorte de él,
    invierte el veredicto: a escala de despliegue la sintonización mejora un 5{,}98\%, con el
    intervalo al 95\% íntegramente por debajo de cero y replicado sobre seis búsquedas
    independientes (Anexo~\ref{anx:tuning_imputador}).

    \item \textbf{En una evaluación por coste, el término de seguridad de la política decide el
    resultado.} Con un término común a todo el catálogo, conocer la demanda real sale más caro
    que la señal censurada sin corregir; con un término proporcional a cada serie, la métrica
    satura hasta no distinguir un imputador razonable de un oráculo. Solo sin término de
    seguridad el coste mide la señal, y la demanda real cuesta cero. El resultado es directamente
    reutilizable por cualquiera que compare señales de demanda por el coste que inducen
    (Sección~\ref{sec:resultados_gemelo}).

    \item \textbf{Medición de la calibración conformal estratificada sobre demanda
    reconstruida.} La variante \textit{Mondrian} se aplica aquí sobre un target reconstruido a
    partir de ventas censuradas, combinación que no cubren los trabajos que llevan la predicción
    conformal al \textit{Newsvendor} ni los que desarrollan la estratificación. Se mide qué le
    ocurre: la cobertura empírica queda por debajo del nivel nominal, la brecha se estrecha al
    ampliar el panel un orden de magnitud sin cerrarse, y el análisis atribuye el resto a la
    no-intercambiabilidad temporal entre calibración y test
    (Sección~\ref{sec:resultados_conformal}).

    \item \textbf{Sistema completo, trazable y reproducible.} Se implementa el recorrido entero
    desde el dato bruto hasta la decisión de pedido: preparación del panel, reconstrucción de la
    censura, ingeniería de características sin fuga temporal, modelos probabilísticos,
    calibración, política de inventario, API operacional, cuadro de mandos y despliegue
    contenerizado con integración continua. Cada ejecución queda registrada con su
    configuración, sus métricas y el \textit{commit} que la produjo, y toda cifra citada en esta
    memoria nombra la corrida que la generó. Esa trazabilidad es la que hace verificables las
    cuatro aportaciones anteriores.
\end{enumerate}

\section{Metodología}
El desarrollo del sistema sigue el ciclo de \textbf{CRISP-DM}~\cite{chapman2000crispdm}
(comprensión del problema, comprensión de los datos, preparación, modelado, evaluación y
despliegue) en su extensión para aprendizaje automático \textbf{CRISP-ML(Q)}~\cite{studer2021crispmlq},
cuya diferencia relevante aquí es que añade una fase final de monitorización y mantenimiento del
modelo desplegado, ausente del estándar original y necesaria porque el riesgo de degradación en un
entorno cambiante es precisamente uno de los fenómenos que este trabajo mide. A ello se suma una
particularidad que condiciona todas las etapas: al tratarse de datos temporales y de decisiones que se toman antes
de conocer la demanda, cada etapa debe respetar el orden del tiempo. Toda variable empleada en una
predicción ha de estar disponible en el instante en que esa predicción se emite, y toda comparación
entre alternativas debe realizarse sobre datos posteriores a los usados para ajustarlas. Este
principio, y no la elección de un algoritmo concreto, es el que gobierna el diseño experimental del
trabajo.

Sobre esa base, el trabajo se organiza en seis etapas encadenadas:

\begin{enumerate}
    \item \textbf{Análisis y preparación de los datos.} Caracterización del conjunto de partida,
    identificación de sus patrones y de sus patologías, y construcción de un panel diario homogéneo
    (Capítulo~\ref{chap:analisis_datos}).
    \item \textbf{Reconstrucción de la demanda censurada.} Estimación de la demanda que se habría
    producido en los días con rotura de stock, mediante tres estrategias alternativas que se
    comparan bajo un mismo protocolo en lugar de adoptarse por defecto
    (Capítulos~\ref{chap:metodologia} y~\ref{chap:resultados}).
    \item \textbf{Ingeniería de características y modelado probabilístico.} Codificación explícita
    de la historia y del calendario en variables sujetas a un contrato causal, y entrenamiento de
    modelos globales que predicen no un valor único sino varios cuantiles de la distribución de
    demanda (Capítulo~\ref{chap:metodologia}).
    \item \textbf{Calibración y validación temporal.} Ajuste de los intervalos de predicción para
    dotarlos de una garantía de cobertura verificable, y evaluación mediante \textit{backtesting}
    con ventanas móviles que reproducen el uso real del sistema
    (Capítulos~\ref{chap:metodologia} y~\ref{chap:resultados}).
    \item \textbf{Traducción de la predicción en decisión.} Conversión de la distribución estimada
    en una cantidad de pedido según la asimetría de costes del negocio, y evaluación del coste
    logístico y del nivel de servicio que esa política induce
    (Capítulos~\ref{chap:metodologia} y~\ref{chap:resultados}).
    \item \textbf{Operativización, monitorización y mantenimiento.} Empaquetado del sistema como
    aplicación desplegable con interfaz de consulta, cuadro de mandos y trazabilidad de cada
    ejecución, de modo que cualquier resultado publicado pueda reproducirse
    (Capítulo~\ref{chap:implementacion}). Esta etapa instancia la fase final de
    CRISP-ML(Q)~\cite{studer2021crispmlq}: detección de deriva sobre las variables de entrada,
    reentrenamiento periódico y una comparación de cadencias de reentreno que cuantifica si
    mantener el modelo al día compensa su coste.
\end{enumerate}

Las etapas no se recorrieron una sola vez ni en línea recta: varias decisiones de diseño se
revisaron al detectar que un resultado provenía de un defecto del protocolo de medida y no del
fenómeno estudiado. El criterio para cerrar cada etapa fue siempre el mismo: una diferencia entre
alternativas solo se considera real si se sostiene sobre datos no empleados para elegirlas y si su
magnitud supera la variabilidad del propio experimento, para lo cual toda comparación entre
alternativas se acompaña de un intervalo de confianza obtenido por remuestreo. El modelo final no
se elige por tanto con una métrica de error estadístico, sino leyendo conjuntamente las dos
magnitudes que la operación percibe, el coste logístico y el nivel de servicio, cada una con su
intervalo; el Capítulo~\ref{chap:resultados} expone qué ocurre cuando ambas no coinciden.

\section{Estructura del Documento}
La presente memoria se organiza en siete capítulos adicionales:
\begin{itemize}
    \item El \textbf{Capítulo 2: Estado del Arte} revisa la literatura actual sobre series temporales, modelos de \textit{boosting}, conformal prediction y gestión de inventarios.
    \item El \textbf{Capítulo 3: Análisis de Datos} describe el dataset utilizado y su preparación.
    \item El \textbf{Capítulo 4: Metodología} detalla la formulación matemática de los modelos, la calibración avanzada, la optimización multi-objetivo y la política de pedido.
    \item El \textbf{Capítulo 5: Implementación y Arquitectura} describe el diseño del software MLOps, la API de despliegue, el sistema de diagnóstico automático y la infraestructura CI/CD.
    \item El \textbf{Capítulo 6: Resultados} presenta el análisis comparativo de los modelos, la evaluación económica de la corrección de censura y la explicabilidad del sistema.
    \item El \textbf{Capítulo 7: Conclusiones} sintetiza los hallazgos principales.
    \item El \textbf{Capítulo 8: Presupuesto} detalla el coste económico del desarrollo del proyecto.
\end{itemize}

La memoria se completa con cuatro anexos. El \textbf{Anexo~\ref{anx:ods}} analiza la
vinculación del trabajo con los Objetivos de Desarrollo Sostenible de la Agenda 2030; el
\textbf{Anexo~\ref{anexo:dashboard}} recoge las capturas del cuadro de mandos web, que
documentan visualmente las vistas descritas en el Capítulo~\ref{chap:implementacion}; el
\textbf{Anexo~\ref{anx:tuning_imputador}} documenta la búsqueda de hiperparámetros del
imputador supervisado, cuyo resultado no modifica el sistema final pero deja una lección
metodológica sobre la validez de este tipo de búsquedas; y el \textbf{Anexo~\ref{anexo:mlflow}}
presenta la trazabilidad de corridas y artefactos en el almacén de experimentos de MLflow.

Por último, dado que el texto emplea terminología específica de la previsión probabilística y de la
gestión de inventarios, se ha incluido al inicio del documento, antes del índice, una
\textbf{Lista de Abreviaturas y Acrónimos} que reúne los términos y siglas utilizados a lo largo de
la memoria. Se recomienda su consulta al lector no familiarizado con el vocabulario del área.

\section{Convenciones}
Con el fin de facilitar la lectura, el texto sigue las siguientes convenciones tipográficas:

\begin{itemize}
    \item Los \textbf{términos técnicos procedentes del inglés} sin traducción asentada se marcan en
    cursiva al introducirse. Los de uso muy frecuente a lo largo de la memoria, como \textit{retail},
    \textit{forecasting} o \textit{stockout}, se emplean después en redonda para no
    recargar la lectura. La cursiva se reserva también para el énfasis puntual.
    \item Los \textbf{elementos de código} (nombres de módulos, ficheros, funciones, columnas de
    datos y claves de configuración) se escriben en tipografía monoespaciada, como en
    \texttt{observed\_demand} o \texttt{configs/experiment.yaml}. Esta tipografía no se emplea para
    ningún otro fin.
    \item Las \textbf{cifras} siguen la convención española: coma como separador decimal y punto como
    separador de millar. Los costes se expresan en \textbf{unidades monetarias} (u.m.) y no en euros,
    ya que el conjunto de datos utilizado no publica precios reales; las comparaciones entre
    políticas son por tanto relativas y no representan importes de un catálogo concreto.
    \item Los \textbf{acrónimos} se desarrollan la primera vez que aparecen, seguidos de la sigla
    entre paréntesis, y se recogen además en la Lista de Abreviaturas y Acrónimos del inicio del
    documento.
    \item Las \textbf{referencias bibliográficas} se citan de forma numerada entre corchetes y se
    detallan en la bibliografía final. Las direcciones web de herramientas, productos o repositorios
    se indican como nota al pie y no como entrada bibliográfica.
\end{itemize}

% !TEX root = ../main.tex
\chapter{Estado del Arte}
\label{chap:estado_del_arte}

Este capítulo presenta una revisión de la literatura y las bases teóricas fundamentales sobre las que se apoya el sistema desarrollado. Se exploran los paradigmas modernos de previsión de demanda, la evolución de los modelos basados en árboles, la cuantificación de la incertidumbre y la optimización de inventarios bajo condiciones de censura.

\section{Previsión de la Demanda en Retail}
Tradicionalmente, la previsión de series temporales se ha abordado mediante modelos \textbf{locales}, donde se entrena un modelo individual para cada SKU o combinación tienda-producto. Los enfoques clásicos como ARIMA \cite{hyndman2021forecasting} o el suavizado exponencial (ETS) han sido el estándar durante décadas gracias a su interpretabilidad y su sólida base estadística. Sin embargo, en el contexto del \textit{retail} masivo presentan tres limitaciones estructurales: (1) su coste computacional crece linealmente con el número de series, haciendo inviable el ajuste individual sobre catálogos de decenas de miles de SKUs; (2) requieren un historial mínimo por serie para estimar sus parámetros, fallando sistemáticamente en productos nuevos o con alta intermitencia (\textit{cold-start problem}); y (3) no pueden aprovechar información compartida entre series como la estacionalidad semanal o la elasticidad de precio, que son patrones transversales a todo el catálogo.

El paradigma actual ha girado hacia los \textbf{Modelos Globales de Pronóstico (GFMs)}, donde un único modelo se entrena sobre la totalidad de las series temporales de la base de datos. Montero-Manso y Hyndman \cite{montero2021principles} demostraron teóricamente que los modelos globales no son restrictivos: bajo condiciones razonables pueden aproximar cualquier función local, pero además se benefician del aprendizaje de representaciones compartidas. Este enfoque ha dominado competiciones de referencia como la M5 Walmart \cite{makridakis2022m5}, donde las soluciones ganadoras fueron casi exclusivamente modelos globales basados en \textit{Machine Learning}.

Una alternativa relevante a los modelos de \textit{boosting} son las arquitecturas de \textit{deep learning} para series temporales, como DeepAR \cite{salinas2020deepar} o el Temporal Fusion Transformer (TFT) \cite{lim2021temporal}. Aunque estas arquitecturas pueden capturar dependencias temporales a largo plazo de forma automática, presentan desventajas en el contexto de este trabajo: requieren volúmenes de datos superiores para converger, son difíciles de interpretar para los responsables logísticos y su tiempo de entrenamiento es sensiblemente mayor sobre datasets de 4,5 millones de observaciones con covariables heterogéneas. En la competición M5 \cite{makridakis2022m5}, los modelos de \textit{boosting} superaron consistentemente a las redes neuronales en términos de precisión y eficiencia, lo que justifica su elección en este trabajo.

\section{Modelos de Gradient Boosting: Comparativa y Selección}
Los algoritmos de \textit{Gradient Boosted Decision Trees} (GBDT) son el estado del arte para datos tabulares y series temporales con covariables exógenas. Las tres implementaciones principales (XGBoost \cite{chen2016xgboost}, LightGBM \cite{ke2017lightgbm} y CatBoost \cite{prokhorenkova2018catboost}) dominaron la competición M5 Walmart y comparten el núcleo algorítmico de boosting por gradiente, pero difieren en aspectos críticos para el contexto de este trabajo.

La Tabla~\ref{tab:gbdt_comparison} resume las dimensiones más relevantes para el problema de previsión de demanda retail con 50.000 series heterogéneas y variables categóricas de alta cardinalidad.

\begin{table}[htbp]
    \centering
    \small
    \setlength{\tabcolsep}{5pt}
    \begin{tabular}{lccc}
        \toprule
        \textbf{Criterio} & \textbf{XGBoost} & \textbf{LightGBM} & \textbf{CatBoost} \\
        \midrule
        Crecimiento de árbol & Por niveles & Por hojas & Oblivious (simétrico) \\
        Variables categóricas & Manual (OHE) & Nativa (particiones) & Nativa (Ordered TS) \\
        Riesgo de \textit{target leakage} & Alto & Medio & Bajo (Ordered Boosting) \\
        Velocidad en datasets grandes & Alta & Muy alta & Media \\
        Regularización implícita & Moderada & Moderada & Alta \\
        Predicción de cuantiles & Sí & Sí & Sí \\
        Incertidumbre nativa (varianza) & No & No & Sí \\
        \bottomrule
    \end{tabular}
    \caption{Comparativa de implementaciones GBDT para el contexto de retail demand fore\-cast\-ing.}
    \label{tab:gbdt_comparison}
\end{table}

La tabla recoge las tres implementaciones que la literatura considera de referencia, pero el
sistema desarrollado implementa \textbf{dos}: LightGBM y CatBoost, ambas bajo el mismo protocolo
de backtesting y el mismo envoltorio cuantílico. XGBoost se descarta antes de implementarse, y por
la fila que más pesa en este dominio: es la única de las tres que exige codificar manualmente unos
identificadores de tienda, producto y categoría de cardinalidad alta, y el \textit{one-hot} de ese
catálogo produce una matriz dispersa que ni escala ni permite al árbol agrupar categorías con
comportamiento parecido.

Entre las dos que sí se implementan, la elección del campeón \textbf{no se decide con los
argumentos de esta tabla}, sino empíricamente y por el criterio que la operación percibe: el coste
logístico simulado sobre el protocolo de backtesting, mediante la lógica de promoción descrita en
la Sección~\ref{sec:promocion}. Esa distinción importa, porque el orden por error puntual y el
orden por coste no coinciden: reportarlos por separado es una de las conclusiones del trabajo. El
campeón resultante es \textbf{CatBoost}, y las propiedades de la tabla explican \emph{por qué es
razonable} que gane, no la razón por la que se selecciona:

\begin{itemize}
    \item \textbf{Gestión nativa de categóricas con estadísticos ordenados:} las dos manejan las
    categóricas sin codificación manual, pero lo hacen de forma distinta. LightGBM ordena las
    categorías por el estadístico del gradiente y parte sobre ese orden, lo que en categorías con
    pocas observaciones puede sobreajustar; CatBoost calcula estadísticos de destino sobre una
    permutación previa (\textit{ordered target statistics}), construcción diseñada precisamente
    para que el valor de la etiqueta de una observación no entre en el estadístico que la codifica.
    Con identificadores de la cardinalidad de este catálogo, esa diferencia es la más relevante de
    las tres.
    \item \textbf{Regularización implícita mediante árboles \textit{oblivious}:} los árboles
    simétricos de CatBoost aplican la misma condición de división en todo un nivel, lo que actúa
    como regularizador estructural sobre un panel donde muchas series tienen historia corta.
\end{itemize}

La estimación nativa de incertidumbre que CatBoost ofrece (última fila de la tabla) \textbf{no se
explota} en este sistema: la cuantificación de la incertidumbre se obtiene de la rejilla de
cuantiles y de la calibración conformal, que son comunes a los dos \textit{backends} y por tanto
mantienen la comparación entre ellos en igualdad de condiciones. Se recoge en la tabla como
propiedad de la herramienta, no como capacidad utilizada.

\section{Cuantificación de la Incertidumbre y Conformal Prediction}
Una predicción puntual es insuficiente para la gestión de inventarios; es necesario cuantificar la probabilidad de diferentes escenarios de demanda para poder aplicar el criterio del Newsvendor. Los enfoques clásicos de cuantificación de incertidumbre se dividen en tres familias: métodos bayesianos, regresión cuantílica y métodos de calibración post-hoc. Los métodos bayesianos \cite{hyndman2021forecasting} ofrecen intervalos bien fundamentados teóricamente pero requieren especificar distribuciones a priori y son computacionalmente costosos a escala. La regresión cuantílica \cite{koenker2005quantile} estima directamente los cuantiles de la distribución condicionada, pero no garantiza cobertura calibrada: el cuantil $q_{0{,}9}$ estimado puede contener la demanda real solo el 70\% de las veces si la distribución de test difiere de la de entrenamiento.

La \textbf{Conformal Prediction (CP)} \cite{angelopoulos2021gentle, vovk2005algorithmic} resuelve este problema ofreciendo garantías de cobertura \textit{distribution-free}: sin asumir ninguna forma funcional de la distribución de la demanda, garantiza que el intervalo predicho al nivel $1-\alpha$ contenga el valor real al menos con probabilidad $1-\alpha$ sobre cualquier muestra intercambiable. Existen varias variantes en la literatura: la CP completa (\textit{full conformal}) es computacionalmente inviable para modelos grandes; la CP dividida (\textit{split conformal}) calibra sobre un conjunto de holdout separado y es eficiente pero desperdicia datos; la CP cruzada (\textit{cross-conformal}) combina las ventajas de ambas mediante validación cruzada.

Un supuesto crítico atraviesa toda la familia: la \textbf{intercambiabilidad} entre los datos de calibración y los de test, que las series temporales violan por construcción. Barber et al.\ \cite{barber2023beyond} caracterizan teóricamente esta situación y acotan la pérdida de cobertura en función de una medida de desviación respecto a la intercambiabilidad, estableciendo que la garantía no desaparece sino que se degrada de forma controlada. Sobre esa base, los métodos de CP adaptativa relajan el supuesto en la práctica: \textit{EnbPI} \cite{xu2021conformal} (Ensemble Batch Prediction Intervals) ajusta dinámicamente el ancho del intervalo basándose en los errores recientes observados, ACI \cite{gibbs2021adaptive} actualiza el nivel de significancia en línea según los fallos de cobertura acumulados, y Zaffran et al.\ \cite{zaffran2022adaptive} comparan sistemáticamente estas variantes sobre series de demanda energética, un dominio con la misma combinación de estacionalidad fuerte y horizonte corto que caracteriza al retail de frescos. Estos métodos son por tanto especialmente robustos ante la no-estacionariedad propia de series de demanda retail con estacionalidad y eventos promocionales; como se verá en el Capítulo~\ref{chap:resultados}, esta limitación es precisamente la que acota la cobertura alcanzada por el sistema desarrollado.

Sin embargo, la CP estándar solo garantiza cobertura \textit{marginal}: el intervalo puede ser correcto en media pero fallar sistemáticamente en subgrupos específicos, como las categorías de alta estacionalidad o los productos intermitentes. La \textbf{Mondrian Conformal Prediction} \cite{vovk2005algorithmic} extiende la garantía a cobertura \textit{condicional} al estratificar la calibración por grupos homogéneos, asegurando que cada subgrupo tenga su propia garantía de cobertura independiente. Esta propiedad es crítica en un catálogo retail heterogéneo.

\section{Optimización de Inventarios: El Problema del Newsvendor}
\label{sec:newsvendor}
La conexión entre el pronóstico y la acción operativa se realiza mediante la teoría de gestión de inventarios estocásticos. El modelo del \textbf{Newsvendor} \cite{silver1998inventory} es el marco de referencia para decisiones de pedido bajo demanda incierta en un único periodo: determina la cantidad óptima a pedir ($Q^*$) minimizando el coste esperado, que equilibra el coste de sobrestock ($c_o$, exceso de producto que se desperdicia o deprecia) y el coste de infra-stock ($c_u$, venta perdida y pérdida de fidelidad del cliente).

La solución óptima viene dada por el \textit{fractil crítico}:
\begin{equation}
    P(D \leq Q^*) = \frac{c_u}{c_u + c_o}
\end{equation}
donde $D$ es la demanda aleatoria. Esta fórmula tiene una interpretación directa: si el coste de rotura es cuatro veces mayor que el de sobrestock ($c_u = 4, c_o = 1$), el fractil óptimo es 0,8, es decir, hay que pedir la cantidad que cubre el 80\% de los escenarios de demanda posibles. Nótese que el fractil óptimo es exactamente el cuantil $q_{0{,}8}$ de la distribución predicha, lo que establece un puente directo entre los intervalos de confianza generados por el modelo probabilístico y la decisión de pedido.

Para entornos de múltiples periodos, el modelo del Newsvendor se extiende a políticas de reposición dinámica. La política $(s, S)$ \cite{silver1998inventory} es la más extendida en la práctica: cuando el nivel de inventario cae por debajo del punto de reorden $s$, se emite un pedido hasta el nivel objetivo $S$. Esta política es óptima bajo costes de pedido fijos y demanda estocástica estacionaria. Su caso límite cuando el coste fijo de pedido es nulo es la política base-stock u Order-Up-To. Este trabajo no modela el estado de inventario entre periodos: resuelve una decisión de un solo periodo por origen de predicción, suficiente para ordenar políticas por su coste y su nivel de servicio, pero no para reproducir la dinámica de reposición que estas familias describen. La extensión a un régimen con estado se recoge como línea de trabajo futuro.

Esta conexión entre predicción y decisión ha cristalizado en la última década en la línea de investigación conocida como \textit{decision-focused learning} o \textit{predict-then-optimize}. Bertsimas y Kallus \cite{bertsimas2020prescriptive} formalizan el paso de la analítica predictiva a la prescriptiva, mostrando que el objetivo relevante no es el error del pronóstico sino el coste de la decisión que induce, y Elmachtoub y Grigas \cite{elmachtoub2022spo} proponen entrenar el modelo directamente contra la pérdida del problema de optimización aguas abajo (\textit{SPO loss}) en lugar de contra una métrica estadística intermedia. El enfoque de este trabajo se sitúa en esta familia pero por una vía más ligera y directamente interpretable: en lugar de diferenciar a través del solver, explota la equivalencia analítica exacta entre el fractil crítico del Newsvendor y el cuantil que minimiza la \textit{pinball loss}, lo que permite alinear entrenamiento y decisión sin abandonar los modelos de \textit{boosting} estándar.

Sin embargo, tanto el Newsvendor como la política $(s, S)$ asumen implícitamente que la distribución $D$ es conocida y libre de sesgo. Esta hipótesis se viola cuando los datos de entrenamiento están censurados por roturas de stock, lo que desplaza el fractil crítico hacia abajo y genera pedidos insuficientes de forma estructural. La sección siguiente analiza este fenómeno y su impacto en la cadena de decisión.

\section{Demanda Censurada y el Efecto de Stockout}
Un desafío crítico y frecuentemente ignorado en \textit{retail} es que las ventas observadas solo coinciden con la demanda real cuando hay stock disponible. Cuando ocurre una rotura de stock (\textit{stockout}), los datos se ven \textbf{censurados por la derecha}: la venta registrada no es la demanda, sino una \emph{cota inferior} de ella, ya que observamos cero (o muy pocas) ventas mientras la demanda latente (lo que los clientes habrían comprado) es estrictamente mayor. Formalmente, al observar $\min(D, \text{stock})$ el valor verdadero queda siempre por encima del dato registrado, que es la definición de censura por la derecha en el sentido del análisis de supervivencia. El estudio teórico de los sistemas de inventario con demanda inobservable se remonta a trabajos fundamentales como el de Nahmias \cite{nahmias1994demand}, y posteriormente, Musalem et al.\ \cite{musalem2010structural} estimaron empíricamente que ignorar este fenómeno subestima la demanda real entre un 5\% y un 30\% dependiendo de la tasa de rotura histórica del producto.

Si un modelo se entrena con estos datos sesgados sin corrección, cae en una ``espiral de muerte'' (\textit{death spiral}): el modelo predice baja demanda, se pide poco producto, se agota antes, y el modelo vuelve a reducir su previsión. Este ciclo se retroalimenta hasta colapsar la reposición de ciertos SKUs. La literatura propone varias estrategias de corrección con distintos compromisos entre precisión y coste computacional:

\begin{itemize}
    \item \textbf{Modelos Tobit}: modelan explícitamente la censura como una distribución truncada, recuperando la media de la distribución latente. Precisos pero requieren especificar la forma distribucional y son computacionalmente costosos a escala.
    \item \textbf{Algoritmos EM (Expectation-Maximization)}: iteran entre estimar la demanda latente y reajustar los parámetros del modelo hasta convergencia. Más flexibles que Tobit pero igualmente costosos para catálogos de 50.000 series.
    \item \textbf{Escalado proporcional}: estima la demanda latente multiplicando la demanda observada por un factor inversamente proporcional a las horas de disponibilidad. Simple y escalable, aunque asume linealidad en la relación entre disponibilidad y demanda.
    \item \textbf{Modelo supervisor}: entrena un modelo auxiliar sobre los días sin stockout para predecir la demanda latente en los días afectados. Captura patrones no lineales y es el enfoque adoptado en este trabajo como estrategia principal.
\end{itemize}

La censura introduce además un segundo problema para la cuantificación de la incertidumbre: los métodos de regresión cuantílica estándar aprenden distribuciones sesgadas si los datos de calibración contienen observaciones censuradas, produciendo intervalos de confianza demasiado estrechos precisamente en los SKUs con mayor historial de roturas. La \textbf{Mondrian Conformal Prediction} mitiga este efecto al calibrar los intervalos de forma separada por categoría de producto: incluso si la distribución global está sesgada, la calibración estratificada corrige el sesgo de cobertura dentro de cada grupo homogéneo, garantizando que el fractil crítico del Newsvendor se calcule sobre una distribución correctamente estimada.

\section{Síntesis y Posicionamiento del Trabajo}
\label{sec:posicionamiento}

La revisión de la literatura revela que los avances en cada uno de los pilares descritos han evolucionado de forma mayoritariamente independiente. Los trabajos sobre modelos globales de \textit{boosting} \cite{makridakis2022m5, ke2017lightgbm, prokhorenkova2018catboost} optimizan la precisión predictiva sin conectarla con la decisión de inventario. Los trabajos sobre Conformal Prediction \cite{angelopoulos2021gentle, xu2021conformal} garantizan cobertura estadística pero asumen datos no censurados y no evalúan el impacto logístico de los intervalos. Los trabajos sobre optimización de inventarios bajo incertidumbre \cite{silver1998inventory, ban2019big} asumen distribuciones de demanda conocidas y libres de sesgo. La literatura de \textit{decision-focused learning} \cite{bertsimas2020prescriptive, elmachtoub2022spo} sí conecta explícitamente predicción y decisión, y es la más próxima al planteamiento de este trabajo, pero lo hace sobre datos limpios y sin cuantificar la incertidumbre con garantías formales ni evaluar la decisión en un régimen dinámico con estado de inventario. Finalmente, la literatura sobre demanda censurada \cite{musalem2010structural} propone correcciones estadísticas pero raramente las integra en un sistema de previsión probabilística con evaluación end-to-end.

Ningún trabajo previo identificado integra simultáneamente estos cuatro elementos: (1) corrección de demanda censurada, (2) forecasting probabilístico global con garantías de cobertura condicional, (3) transformación de la distribución predicha en una decisión de inventario vía el fractil crítico del Newsvendor, y (4) evaluación del impacto económico de esa decisión en las dos magnitudes que la operación percibe, coste logístico y nivel de servicio. Esta integración constituye la contribución central del presente trabajo y justifica el diseño del sistema descrito en los capítulos siguientes.

% !TEX root = ../main.tex
\chapter{Análisis del Conjunto de Datos}
\label{chap:analisis_datos}

Antes de abordar el modelado, es imperativo comprender la naturaleza y estructura de los datos. Este capítulo presenta un Análisis Exploratorio de Datos (EDA) sobre el dataset \textit{Fresh RetailNet}, identificando los patrones, retos y oportunidades que definen el problema de previsión.

\section{Descripción del Dataset}
El conjunto de datos \textit{Fresh RetailNet} es un benchmark masivo diseñado específicamente para el sector retail. Contiene series temporales diarias de ventas de productos frescos en diversas tiendas, acompañadas de metadatos jerárquicos y variables exógenas.

Las características principales del dataset son:
\begin{itemize}
    \item \textbf{Panel de Datos}: Series temporales de longitud variable (hasta 90 días) para miles de combinaciones tienda-producto (\textit{SKUs}).
    \item \textbf{Variables de Objetivo}: Cantidad vendida (\textit{sale\_amount}) normalizada.
    \item \textbf{Variables de Contexto}: Precios, descuentos, festivos y datos meteorológicos (temperatura, precipitación).
    \item \textbf{Indicadores de Stockout}: Información crítica sobre las horas en las que el producto no estuvo disponible en el estante.
\end{itemize}

\section{Patrones Temporales y Estacionalidad}
El análisis visual de las series revela una fuerte estacionalidad semanal. Las ventas tienden a concentrarse en los días previos al fin de semana, con caídas significativas en días festivos oficiales cuando las tiendas permanecen cerradas. Esta naturaleza cíclica justifica el uso de \textit{lags} de 7 días y variables de calendario en la fase de ingeniería de características.

\section{Análisis de la Intermitencia}
Un rasgo distintivo del \textit{retail} fresco es la intermitencia. Muchas series presentan días con ventas cero, no debido a una falta de interés del consumidor, sino a la falta de disponibilidad del producto. Se ha cuantificado que aproximadamente el 15\% de las observaciones del dataset están afectadas por algún grado de \textit{stockout}, lo que subraya la importancia de la reconstrucción de la demanda latente propuesta en este trabajo.

\section{Correlación con Variables Exógenas}
Se observa una correlación positiva moderada entre la temperatura media y la demanda de ciertas categorías de productos frescos (ej. frutas de temporada). Asimismo, el factor "descuento" actúa como un fuerte impulsor de la demanda, alterando significativamente el patrón basal de ventas. El modelo desarrollado debe ser capaz de integrar estas señales para superar a los modelos autorregresivos simples.

\section{Interpretación de las Visualizaciones del EDA}

Las visualizaciones generadas automáticamente complementan el análisis tabular del panel preparado y permiten identificar de forma más clara la estructura temporal, la heterogeneidad entre series y el impacto operativo de los \textit{stockouts}. Su objetivo no es modificar la semántica del problema, sino aportar evidencia descriptiva que justifique las decisiones posteriores de modelado y validación.

\subsection{Cobertura temporal y calidad estructural del panel}

La figura de cobertura temporal (\texttt{coverage\_heatmap.png}) muestra la presencia de observaciones para cada serie a lo largo del intervalo analizado. Su principal utilidad es verificar si el panel final presenta huecos, discontinuidades o series truncadas. En este caso, la cobertura es homogénea, lo que respalda que el conjunto preparado utilizado en el modelado mantiene una estructura balanceada tras el preprocesamiento y el filtrado de series.

\subsection{Distribución de la demanda observada}

La distribución global de la demanda (\texttt{observed\_demand\_distribution.png}) evidencia que la variable objetivo no sigue una forma aproximadamente normal, sino que presenta concentración en rangos bajos y una cola hacia valores mayores. Esto es coherente con un contexto \textit{retail} de producto fresco, donde la demanda diaria suele ser moderada, irregular y sensible a factores operativos y contextuales. La visualización de las series de mayor volumen mediante diagramas de caja (\texttt{observed\_demand\_boxplot\_top\_series.png}) confirma, además, que existe heterogeneidad relevante entre series, tanto en nivel medio como en variabilidad. Esta evidencia apoya el uso de modelos globales con contexto, en lugar de asumir que todas las series comparten una dinámica homogénea.

\subsection{Estacionalidad semanal}

La figura \texttt{weekday\_demand\_profile.png} muestra la evolución de la demanda media y mediana según el día de la semana. El patrón observado sugiere una estacionalidad semanal clara, con diferencias sistemáticas entre días laborables y fin de semana. Esta regularidad justifica directamente la inclusión de variables de calendario y de retardos semanales en la etapa de ingeniería de características. Desde el punto de vista metodológico, esta es una de las señales más sólidas del EDA, ya que conecta de forma clara la estructura de los datos con el diseño de los predictores temporales.

\subsection{Intermitencia y demanda cero}

La figura \texttt{zero\_demand\_rate\_by\_series.png} permite identificar las series con mayor proporción de observaciones en cero. Su lectura muestra que la intermitencia no afecta por igual a todo el panel, sino que se concentra con distinta intensidad según la serie. Esta heterogeneidad es importante porque advierte que las métricas agregadas pueden ocultar comportamientos muy diferentes entre productos o combinaciones tienda-producto. En consecuencia, la dificultad del problema no debe entenderse como uniforme, sino como dependiente del régimen operativo y de la naturaleza concreta de cada serie.

\subsection{Frecuencia e intensidad de los stockouts}

La distribución de horas de \textit{stockout} (\texttt{stockout\_hours\_distribution.png}) muestra que la falta de disponibilidad no es un fenómeno anecdótico, sino una característica estructural del conjunto analizado. Esto refuerza la idea de que el problema de predicción no puede interpretarse como una simple serie temporal de ventas observadas, ya que parte de esas observaciones están condicionadas por restricciones de disponibilidad en el lineal.

La figura \texttt{stockout\_band\_demand.png} profundiza en este punto al comparar la demanda observada entre bandas de intensidad de \textit{stockout}. El patrón es coherente con la hipótesis de censura: cuando las horas de \textit{stockout} son severas, la demanda observada tiende a reducirse. No debe interpretarse esto como una estimación exacta de demanda latente, pero sí como evidencia descriptiva de que las ventas registradas pueden infraestimar el interés real del consumidor bajo condiciones de falta de stock.

\subsection{Relación entre stockout y demanda observada}

La figura \texttt{stockout\_vs\_demand\_scatter.png} representa la relación entre horas de \textit{stockout} y demanda observada, junto con una tendencia promedio. La relación agregada es negativa, aunque con alta dispersión. Esto indica que el efecto no es determinista ni uniforme, pero sí suficientemente visible como para considerar \texttt{stockout\_hours} una variable contextual relevante. La interpretación correcta es descriptiva y no causal: la visualización no demuestra por sí sola el mecanismo exacto de pérdida de ventas, pero sí aporta evidencia compatible con la existencia de censura operativa.

\subsection{Relaciones con covariables exógenas}

El mapa de correlaciones (\texttt{correlation\_heatmap.png}) y la cuadrícula de relaciones entre covariables y demanda (\texttt{covariate\_vs\_demand\_grid.png}) permiten estudiar la asociación entre la demanda observada y variables como descuento, temperatura o precipitación. En general, las relaciones marginales aparecen como moderadas o débiles, lo que sugiere que estas señales exógenas no deben interpretarse de forma aislada. Sin embargo, su inclusión sigue siendo razonable, ya que un modelo flexible puede capturar interacciones no lineales y dependencias condicionadas por serie o régimen que no se aprecian en una correlación global simple.

\subsection{Concentración del volumen y series representativas}

La figura \texttt{top\_series\_total\_demand.png} muestra que parte del volumen agregado se concentra en un subconjunto reducido de series. Esto contextualiza el uso de filtros como \texttt{top\_n\_series} y explica por qué el comportamiento agregado del panel puede cambiar cuando se restringe el análisis a las series más relevantes por volumen.

Finalmente, la figura \texttt{representative\_series\_panels.png} ofrece una visión cualitativa especialmente útil para interpretar el problema. En estas series se observan simultáneamente patrones semanales, cambios de escala entre productos y episodios donde la demanda observada coexiste con intensidades elevadas de \textit{stockout}. Esta visualización es valiosa porque resume de forma intuitiva la complejidad conjunta del problema: estacionalidad, heterogeneidad entre series y posible censura por falta de disponibilidad.

\subsection{Conclusiones del EDA}

En conjunto, las visualizaciones permiten extraer varias conclusiones relevantes para el diseño experimental del sistema. En primer lugar, el panel preparado presenta una estructura temporal consistente y adecuada para validación temporal. En segundo lugar, existe una estacionalidad semanal clara que justifica el uso de retardos y variables de calendario. En tercer lugar, la heterogeneidad entre series y la intermitencia parcial del panel aconsejan evitar simplificaciones excesivas en el modelado. Por último, la frecuencia y la intensidad de los \textit{stockouts}, junto con su asociación negativa con la demanda observada, refuerzan la idea de que el contexto operativo debe incorporarse explícitamente en el sistema de predicción y en la interpretación de sus resultados.

\section{Visualizaciones Generadas Automáticamente}
\IfFileExists{figures/eda/eda_figures.tex}{
    \input{figures/eda/eda_figures.tex}
}{
    % Ejecuta el modulo de EDA para materializar estas figuras.
}


% --- Capítulos pendientes (estructura) ---
\chapter{Metodología}
\textit{En desarrollo. Este capítulo detalla la formulación matemática de los modelos probabilísticos (Mondrian Conformal Prediction, fractil crítico del Newsvendor), la corrección de demanda censurada, la optimización multi-objetivo NSGA-II y el diseño del gemelo digital logístico $(s,S)$.}

\chapter{Implementación y Arquitectura}
\textit{En desarrollo. Este capítulo describe el pipeline de datos, el diseño MLOps (MLflow, FastAPI, CI/CD), la lógica de promoción del modelo campeón, el diagnóstico automático post-mortem y el cuadro de mandos React.}

\chapter{Resultados y Evaluación Logística}
\textit{En desarrollo. Este capítulo presenta el backtesting walk-forward ($K=3$, $H=7$ días), la calibración conformal, el impacto económico en el gemelo digital, la optimización bajo restricciones de capacidad y la explicabilidad SHAP.}

\chapter{Conclusiones y Trabajo Futuro}
\textit{En desarrollo. Este capítulo sintetiza las contribuciones principales, discute las implicaciones para el sector retail y propone líneas de investigación futuras.}

\chapter{Presupuesto del Proyecto}
\textit{En desarrollo. Este capítulo detalla el coste económico del proyecto: personal, recursos materiales e infraestructura, gastos generales y presupuesto total con IVA.}

\backmatter
\bibliographystyle{unsrt}
\bibliography{references}

\end{document}
